% !TeX root = ../main.tex

\documentclass[../main.tex]{subfiles}

\begin{document}

Phụ lục này trình bày các đoạn mã Python được sử dụng để nhận dạng tham số và thiết kế bộ điều khiển cho hệ thống điều khiển cascade.

\section{Mã nguồn nhận dạng đối tượng quán tính bậc nhất có trễ (Q1T)}

Mã nguồn nhận dạng tham số cho mô hình Q1T $O_1(s) = \frac{K_v}{T_v s + 1}e^{-\tau_v s}$ của đối tượng van điều khiển:

\lstinputlisting[language=Python, style=tech, caption={Mã nguồn nhận dạng Q1T}, label={code:q1t}, inputencoding=utf8, extendedchars=true]{../qtb1t.py}

\section{Mã nguồn nhận dạng đối tượng quán tính bậc hai có trễ (Q2T)}

Mã nguồn nhận dạng tham số cho mô hình Q2T $O_2(s) = \frac{K_{bqn}}{(T_1 s + 1)(T_2 s + 1)}e^{-\tau_{bqn} s}$ của đối tượng bộ quá nhiệt:

\lstinputlisting[language=Python, style=tech, caption={Mã nguồn nhận dạng Q2T}, label={code:q2t}, inputencoding=utf8, extendedchars=true]{../qtb2t.py}

\section{Mã nguồn nhận dạng đối tượng tương đương và thiết kế bộ điều khiển vòng ngoài}

Mã nguồn nhận dạng đối tượng tương đương $G(s) = W_2(s) \cdot O_1(s)$ dạng Q1T và tính toán tham số bộ điều khiển PID vòng ngoài theo Ziegler-Nichols:
\begin{align}
K_{p2} &= \frac{1.2 T_e}{K_e \tau_e}, \quad T_{i2} = 2 \tau_e, \quad T_{d2} = 0.5 \tau_e
\end{align}

\lstinputlisting[language=Python, style=tech, caption={Mã nguồn nhận dạng đối tượng tương đương và thiết kế PID vòng ngoài (Ziegler-Nichols)}, label={code:draw}, inputencoding=utf8, extendedchars=true]{../draw.py}

\section{Mã nguồn tính toán đối tượng tương đương vòng ngoài}

Mã nguồn tính toán và mô phỏng đáp ứng bậc thang của đối tượng tương đương vòng ngoài $V_1(s) = W_2(s) \cdot O_1(s)$ dạng Q2T với $\tau_{V1} = \tau_v + \tau_{bqn}$:

\lstinputlisting[language=Python, style=tech, caption={Mã nguồn tính toán đối tượng tương đương vòng ngoài $V_1(s)$}, label={code:draw2}, inputencoding=utf8, extendedchars=true]{../draw2.py}

\section{Giải thích các thành phần chính}

Các đoạn mã tuân theo cấu trúc: đọc dữ liệu từ CSV, định nghĩa hàm mô hình, xây dựng hàm mục tiêu (tổng bình phương sai số), thiết lập giá trị khởi tạo và ràng buộc, tối ưu hóa bằng \texttt{scipy.optimize.minimize()}, đánh giá kết quả qua $R^2$ và vẽ đồ thị so sánh.

\end{document}